\chapter{Implementation}

This chapter describes the implementation of the major modules of the \textbf{CricketVision: AI-Powered Cricket Performance Analytics and Biomechanics Platform}. The application was implemented using React 18, Tailwind CSS, Node.js, Express.js, MongoDB with Mongoose ODM, Recharts, HTML5 Canvas API, and supporting libraries for document processing, authentication, graphics rendering, and AI-assisted strategy advice.

\section{User Authentication}
The authentication module provides registration and login functionality supporting persona presets (\texttt{Rahul Dravid} - Coach, \texttt{Virat Kohli} - Batter, \texttt{Jasprit Bumrah} - Bowler, \texttt{Alex Morgan} - Analyst) and custom user accounts.

During registration, the user's password is processed using bcrypt password hashing before it is stored in the database. The password is therefore stored as a hash rather than being stored directly.

During login, the entered password is compared with the stored password hash. If the credentials are valid, the user is allowed to access the application and a session context with role permissions is established.

\section{Roster Creation and Player Seeding}
The roster management module allows users to enter, edit, save, and manage player performance records through structured forms.

The form supports information such as:
\begin{itemize}
    \item Personal details (display name, country/franchise, primary role).
    \item Physical workload indicators (Fatigue Level Index, Clutch Rating).
    \item Radar skill attributes (Power, Consistency, Spin, Pace, Fielding, Clutch).
    \item Phase-wise performance metrics (Powerplay, Middle, Death overs).
    \item Biomechanical joint flexion summaries ($\theta_{\text{elbow}}$, shoulder tilt, action status).
\end{itemize}

On server cold boot, the backend controller checks the MongoDB player collection size. If fewer than 100 documents exist, an automated seeding script executes, inserting 104 verified international and IPL player profiles into the database.

\section{360° Wagon Wheel SVG \& Pitch Heatmaps}
The visual graphics module \texttt{PitchAndWagonWheel.jsx} converts polar shot coordinates $(r, \theta)$ into Cartesian 2D SVG canvas points $(x, y)$ using the mathematical transformation:

\begin{equation}
x = x_{\text{center}} + r \cdot \cos\left(\theta - \frac{\pi}{2}\right), \quad y = y_{\text{center}} + r \cdot \sin\left(\theta - \frac{\pi}{2}\right)
\end{equation">

The component renders interactive stadium sector overlays across 8 field sectors (Offside, Legside, Cover, Mid-Wicket, Fine Leg, Third Man, Long On, Long Off) alongside pitch length distribution heatmaps (Full, Good, Short, Yorker).

\section{Stochastic Match Simulation Engine}
The match simulation engine evaluates live target scores, overs elapsed, wickets lost, pitch friction factors, and weather parameters to calculate dynamic win-probability curves $P_{\text{win}}$ recalculated ball-by-ball:

\begin{equation}
P_{\text{win}} = \frac{1}{1 + e^{-\left( \alpha \cdot (RRR_{\text{target}} - RRR_{\text{current}}) + \beta \cdot W_{\text{rem}} + \gamma \cdot F_{\text{pitch}} + \delta \cdot \Omega_{\text{weather}} \right)}}
\end{equation">

The output is rendered dynamically on the frontend using interactive Recharts line graphs.

\section{Pose Estimation Video Biomechanics Analyzer}
The video biomechanics analyzer module loads bowling action video feeds onto an HTML5 canvas element. The keypoint extraction engine calculates the elbow joint flex angle $\theta_{\text{elbow}}$ using the vector dot-product formula:

\begin{equation}
\theta_{\text{elbow}} = \arccos\left( \frac{\vec{BA} \cdot \vec{BC}}{\|\vec{BA}\| \|\vec{BC}\|} \right) \times \left(\frac{180}{\pi}\right)
\end{equation">

If $\theta_{\text{elbow}}$ exceeds the ICC 15-degree flexion threshold, a red warning alert (\textbf{ACTION ILLEGAL: FLEXION $> 15^{\circ}$}) is overlaid on the video canvas.

\section{AI-Assisted Strategy Advice}
The AI Coach module connects to the Google Gemini API to analyze squad stats, fatigue alerts, and upcoming match venue conditions, generating natural language tactical recommendations for team selection, bowling rotations, and rehabilitation schedules.

\section{Team Builder XI Balance Optimizer}
The \texttt{TeamBuilderView.jsx} component evaluates selected 11-player lineups and computes an overall Team Balance Score based on batting depth, pace/spin bowling distribution, and clutch experience ratings.

\section{Database Integration}
MongoDB with Mongoose ODM is used for storing application data.

The backend communicates with the database to create, store, retrieve, update, and prune user and player records. The stored information is reused by the visual analytics, match simulator, pose analyzer, and AI Coach modules.

\section{Frontend and Backend Communication}
The React frontend communicates with the Node.js and Express.js backend through asynchronous HTTP REST API requests.

The backend receives API requests, performs required processing, interacts with the database or AI services, and returns JSON payloads to the frontend.

\section{Integration of Modules}
The completed application integrates the different modules into a common workflow. The user can log in, manage squad rosters, view 360-degree SVG wagon wheels, simulate ball-by-ball match win probabilities, audit video bowling kinematics, and receive AI Coach strategy advice.

The use of structured player documents throughout these modules allows the same data to be reused for multiple outputs.
